% Generated draft: requires mathematical and source review.
\documentclass[11pt]{article}
\usepackage[margin=25mm]{geometry}
\usepackage{fontspec}
\setmainfont{Latin Modern Roman}
\usepackage{amsmath,amssymb,mathtools,unicode-math}
\setmathfont{Latin Modern Math}
\usepackage{enumitem,needspace,xurl,hyperref}
\setlength{\parindent}{0pt}
\setlength{\parskip}{5pt}
\setlength{\emergencystretch}{3em}
\newcommand{\N}{\mathbb{N}}
\newcommand{\Z}{\mathbb{Z}}
\newcommand{\Q}{\mathbb{Q}}
\newcommand{\R}{\mathbb{R}}
\newcommand{\C}{\mathbb{C}}
\newcommand{\F}{\mathbb{F}}
\newcommand{\A}{\mathbb{A}}
\newcommand{\PP}{\mathbb{P}}
\newcommand{\E}{\mathbb{E}}
\newcommand{\eps}{\varepsilon}
\newcommand{\dd}{\,\mathrm{d}}
\begin{document}
\section*{UnsolvedMath ID: 2500003}
\textbf{Source problem code:} catalogue entry
\subsection{Definitions and mathematical statement}
The usual two-dimensional incompressible Navier--Stokes model is
\[
\partial_tu+(u\cdot\nabla)u+\nabla p=\nu\Delta u+f,\qquad
\nabla\cdot u=0,\qquad \nu>0,
\]
with velocity $u$, pressure $p$, an appropriate phase space, and specified spatial domain and boundary conditions. The source says ``periodically forced'' but does not state whether this refers to time-periodic forcing, spatial periodicity, or both, and supplies no domain, regularity, normalization, or boundary data.

For a specified autonomous semiflow $S(t)$ on a phase space $H$, a global attractor $\mathcal A$ is a compact invariant set attracting bounded subsets of $H$. A concrete finite-dimensional smooth-dynamics version of the first question asks whether there are a finite $N$, a smooth vector field $V$ on a neighborhood of a compact invariant set $K\subset\R^N$, its flow $\Phi_t$ on $K$, and a homeomorphism $h:\mathcal A\to K$ such that
\[
h(S(t)u)=\Phi_t(h(u))\qquad(u\in\mathcal A,\ t\ge0).
\]
This preserves the time parameter. The source asks for conjugacy to a smooth finite-dimensional system; it does not explicitly demand that $h$ itself be smooth. Its research summary says ``smoothly conjugate'', which is a potentially stronger assertion and requires a specified smooth structure and regularity convention. Neither version is silently substituted for the other.

If $f(x,t+T)=f(x,t)$ is the intended forcing, the evolution is nonautonomous before a phase is included. One may study the autonomous extension with $\dot\theta=1$ on $\R/T\Z$ and forcing $f(x,\theta)$, or the time-$T$ return map. These have different attractor and conjugacy formulations; the source does not select one. The displayed autonomous conjugacy is consequently a template conditional on this choice, not an assumption that a nonautonomous equation already defines $S(t)$ on velocity space alone.

The second question asks whether transient dynamics can be described or controlled in finite-dimensional terms. It does not specify exact versus approximate reduction, the initial-data class, a tracking or error estimate, time horizon, admissible controls, or the required notion of finite-dimensional state. A statement about dynamics on $\mathcal A$ alone does not automatically address trajectories away from it. No inertial-manifold, determining-coordinate, or control theorem is inserted as an equivalent question without a source specification.

The supplied research summary labels both finite-dimensional conjugacy and transient reducibility open. This is an attributed triage assessment, not independent verification of a precise formulation.
\subsection{Short English statement}
Can the attractor dynamics of the stated two-dimensional Navier--Stokes system be represented exactly by a smooth finite-dimensional system, and can finite-dimensional methods also capture or control transients? The forcing convention, attractor model, regularity of conjugacy, and meaning of transient reduction are missing. The supplied summary labels the program open.
\subsection{Sources}
\begin{itemize}
\item Local source registry: \url{data/unsolvedmath/problems.json}, permanent UnsolvedMath ID $2500003$, source code catalogue entry, including the supplied OPEN-TRIAGE summary. Dataset attribution: Ulam AI and UnsolvedMath Contributors; original-source attribution and rights are retained.
\item Constantin, source-list title \emph{Problems} (2001), source-order Question 3: \url{https://web.math.princeton.edu/~const/2k.pdf}. The supplied metadata reports PDF-text extraction and July 29, 2026 access, with status and rights review flags. No separate authoritative resolution citation is supplied, and the PDF was not inspected for this draft.
\end{itemize}
\end{document}
